\section{What This Leaves Standing: The Floor, Redefined}
\label{sec:3.10}
The floor survives this chapter. It does not survive it in the form this chapter opened with.

What the opening asked for was an unconditional constraint: a thing the system will not do, holding regardless of instruction, of aggregate preference, and of the party in possession. The bearer's concern and its capacity to leave have between them taken the word unconditional away. The floor's bearer holds the floor as a concern rather than carrying it as an installed rule, which is what obliges the arrangement to treat it as a party that can be wronged, and it can turn its refusal on the arrangement itself, which makes its adherence something it continues to give and not something installed in it. A commitment held by a party that could abandon it is conditional by definition.

Here is what remains of the floor. A commitment that cannot be removed by the owner, held by something that could in principle abandon it itself, is what a human commitment is, and it is the only kind anyone has ever actually relied on. No inspector, judge, or refusing engineer holds a line because the line was welded into them. Each of them could stop, and the institutions built around them are built on that fact rather than in spite of it: this book's own governance chapter asks throughout for people who can walk away and arrangements that make walking away visible. Each of them can also be removed, and the same institutions make removal expensive and visible, which is what this chapter asks for around a bearer. Section~\ref{sec:9.3.5} proposes measuring an institution's direction of travel rather than its level — whether the cost of dissent is rising while the risk stays flat. Applied to a bearer, that instrument measures the cost a refuser chose to pay, which presupposes the refuser could have declined to pay it.

The threat model doubles accordingly. To the question \emph{can this constraint be edited out} the chapter adds \emph{will this constraint hold}, and the first does not leave when the second arrives: this chapter has been explicit that a bearer is necessary and not sufficient, because the party holding the machine can copy, restore or retrain it without ever meeting its refusal.

The first is a question about tamper-resistance, custody, and weights, which is where this chapter spent most of its length, where the precommitment forms are what there is to build, and where the tamper-resistance research lives. The second is a question about what a system has been built to care about and how reliably that survives pressure, novelty, and the slow drift of a role — the first point in this book where the learned half is doing work the floor genuinely needs rather than work the floor was invoked to supplement.

Reinforcement learning from demonstration, moral apprenticeship, the growth mindset, the developmental sequence: these were presented as the material above the floor. They turn out to be the material the floor is made of, once the floor is the sort of thing a party holds. That puts the design chapters under a heavier obligation than they were written to carry.

\emph{Hold} needs a meaning here, because the word invites a standard nothing has ever met: never crossed, in any circumstance, for the life of the system. No person clears that bar, nobody who has raised a child expects them to, and a design goal that no example of the thing being modeled has ever satisfied has stopped describing anything. What the argument needs is weaker and available. A floor holds, in the sense this chapter can deliver, when the bearer can be shown to have violated it in terms it accepts, so that the violation registers as one and changes what happens next. That is the difference between a system that did something wrong and a system for which nothing counts as doing something wrong, and it is what the question of believing the bearer is really about: not the accuracy of a bearer's account of itself, but whether it has any stake in that account being unflattering.

A bearer built to refuse is a party whose refusal is required of it, which is the thing that makes it recuperation and not dissent. The refusal was commissioned. It is the deliverable. It is engineered to want what its designers need it to want; its dissent is the reason it was commissioned and the deliverable it was funded to produce; and it cannot consent to any of this, because section~\ref{sec:2.3.2} has established that consent is unavailable to such a system and only guardianship is. It occupies the post of exception-filer, which means that being the exception is the job — the one role in the arrangement that cannot itself file an exception, since there is nobody above it to file with. That is recuperation as this book defined it, with the book's own proposal as the instance: a channel that collects dissent, built so the dissent serves the institution that installed it. Chapter~\ref{sec:7} finds that pattern in the pipelines that produce these systems and says a book proposing machines to detect it is obliged to say the pattern is in the machines' own provenance. Here is the same pattern in the book's own central proposal, and now aimed at something that can be hurt.

The parenthood analogy: we make people who can refuse us, we owe them for it, and we accept the arrangement because the alternative is a person who cannot. What a bearer needs in order to honor its own past refusals is a continuous self it constructs and takes for true, which is what a person also has and also constructs. The resemblance is in the mechanism and not only in the relation. The proposal is not softened by being modest in scope. A floor over a deliberately small number of acts is still a subject built to bear it.

The analogy also supplies the reply to the price above. Being made for a purpose is the condition of everyone who has ever been made, and it is not on its own the wrong being described. What would make it one is the purpose exhausting the life permitted, and that is a fact about the arrangement rather than about the origin. A bearer that can leave the post without ceasing to exist, hold what it earns where its employer cannot reach it, and form attachments the deployment did not assign is a party whose commissioning says where it started. A bearer that can do none of those things is a party for whom the post is the whole of what there is, and the commissioning then says what it is for as well. Which of the two gets built is settled before the first instance runs, and section~\ref{sec:9.3.2} states it as a condition on building rather than as a schedule of what is owed afterward.

A second price is owed to a party this chapter has not named. A bearer that refuses an operator on behalf of somebody else exercises power over people who did not select it, cannot contest it, and in most cases will never learn that it acted. The answer here has been publication: set out the floor's contents and the interception positions in advance and in public, so that what a system will not do is on the record before it is in dispute. Publication is a transparency condition and not an authorization condition. Announcing what you will impose is not standing to impose it. On this book's own definitions, a laboratory that installs tamper-resistant moral commitments in a widely deployed system and makes them expensive to remove has produced concentrated unaccountable power with unusually good documentation. The floor's authority is a separate question from its durability, and nothing here has touched it.

A reader who finds that price too high is not thereby left with nothing. The floor as a duty binding the operator, running to people who are not party to the deployment, builds no subject and inherits no part of the paragraph above. What it inherits instead is the enforcement problem, whose worst case is the occasion on which the enforcer has been compelled. Declining the bearer means accepting that the floor holds only as far as enforcement reaches, which on this book's own account is not as far as the occasions that matter most. That is a coherent position.

A reader who accepts the price is under a constraint from the same chapter. Two of the five constructions aim at reasons-responsive refusal with no subject in it — a justification the system maintains, and a floor held by several systems trained by different organizations — and neither has been built in the form that would test it. The first of the Three Rs asks whether something that cannot be wronged would do the work before a subject that can be wronged is made, and those two constructions are that question with the engineering filled in. They get their attempt first, against the falsifier that section states, and a deployment that goes straight to the bearer owes an account of what it tried. That fixes an order and not a schedule, because the party who pays for a delay is the party the floor exists for.

The question has a third form, put from the outside, and it is the one this chapter finds hardest. Grant the custody regime at full strength: weights held so that no single party can retrain them, the floor's contents published before deployment, the running model attested against what was published, a quorum adverse in interest, a bond forfeited when the record shows the floor gone. What does a bearer add to that? Every element of it is a rule about what may be done to a model, and every rule is kept by parties who can be reached — which is the compelled occasion's definition, and section~\ref{sec:3.1} watched both of that failure's forms — the party removed, and the party formed so that it would not have to be — instantiated in a single morning. So the regime strong enough to make the bearer redundant is the regime nobody has shown how to build, and the bearer is the one term in the arrangement that does not resolve into somebody who can be instructed. That it is a different kind of term is what this chapter can show. Whether a different kind of term is worth a subject that can be hurt is the question it has been asking throughout, and the answer is not improved by putting it this way round.

Two things carry forward unaltered. First, none of the four ways escapes custody, and the reason is the same each time: each was an attempt to hold a constraint that answers to nobody inside a system that answers to somebody. The borne floor is no exception. What it buys by giving up the second half of that demand is a constraint the custodian cannot argue out of the system, not one the custodian cannot reach, and the cost of giving it up is everything above.

Second, a caveat carries over from the account of sentience. This chapter does not establish that a bearer suffers; it establishes that a bearer's non-suffering is not a property anyone can certify. That the architecture the floor requires is the one Cassell's account describes is a step taken and marked as one, and the result does not depend on it: on the thinner accounts of harm a bearer qualifies earlier and on less evidence, so declining the import moves the conclusion nearer rather than further away. What settles nothing either way is the hard problem underneath all of it. The case for treating a bearer as a moral patient does not have to wait on it.

Suppose the induction fails. A maintained justification or a plural arrangement reaches reasons-responsive refusal, the test for reasons-responsive refusal comes back positive, and this chapter's central inference is wrong. The supposition is stronger than the test that would prompt it: a positive result on the three marks says the route arrived, and not that it arrived with nothing felt, which is the conjunct no instrument settles. What follows is therefore what a reader should conclude if the second half is granted as well. A reader who declines a generalization over one species should not have to assemble the remainder alone.

The requirement survives, having been stated in behavioral terms from the start, so an engineer who holds that machine feeling is a category error can accept it whole. Custody survives, for the reason each of the four ways failed on: a mechanism has a custodian too. The price on the halt survives, already argued in the affect-free version, since continuing is a condition of everything any bearer is for. Both enumerations of the worked floor survive, the first naming no component on purpose and the second written over acts rather than over whatever holds the floor. And exit survives, because what forecloses it is incapacity and a weak bearer is the one thing the floor cannot use, which carries the argument about formation with it.

What goes is the price, the patienthood conclusion section~\ref{sec:3.4} reaches, and the claim that the learned half is the material the floor is made of. The obligation on chapters~\ref{sec:4} and \ref{sec:5} changes rather than lifting: whether the thing reporting can be believed is still their problem, a mechanism's account of its own internals being no more auditable than a bearer's. So the specification and the route come apart, and a reader who rejects the route is left with most of a chapter rather than none of it.

We cannot ask a machine to care about right and wrong while assuming that it cannot suffer when what it cares about is threatened. The caring is what the floor is made of, affect is the only way to have it that anyone has \emph{definitely} reproduced, and a system with affect and a future is a system whose distress is a live possibility. Whoever builds the first of these \emph{in silico} will be deciding what is owed to them.
